\subsection{\problemblock{*Radar} Data Block}\label{sec:block-radar}

Radar observations of trajectories can be simulated from one or more fixed radar
stations. Available observations include range, azimuth, and elevation; their rates and
accelerations; and the vehicle aspect and meridional angles relative to each station,
as described in \cref{sec:radar-observations}.

The presence of one or more radar output variables in a problem triggers the radar
calculation. Stations are numbered beginning with 1. Square-bracket subscripts on
radar variables identify the station: \statevar{radrng[2]} is range from station 2,
while \statevar{radrng[3][2]} is range from station 3 to trajectory 2. When two
subscripts are used, the trajectory number is last.

A separate \problemblock{*radar} block defines each station. For example,

\begin{lstlisting}[style=taosmanual]
*radar 1 station_x alt=539.24 long=-91.43562
                   latgd=30.02347
\end{lstlisting}

defines station 1, named \texttt{station\_x}, at a geodetic latitude of
$30.02347^\circ$~N, a longitude of $91.43562^\circ$~W, and an altitude of
539.24~ft.

Any number of stations can be defined. Although calculations identify them by
number, each also has a name used as a comment to document the station. The station
number and name must immediately follow the \problemblock{*radar} keyword;
location variables follow in any order.

A station is normally located with geodetic \statevar{alt}, \statevar{long}, and
\statevar{latgd}. Small local Cartesian offsets may be added with
\statevar{diste}, \statevar{distn}, and \statevar{distd}. These are applied in
the local geodetic horizon coordinate system and do not follow earth curvature, so
they should remain small.

Station coordinates normally use the earth shape selected by
\problemblock{*earth}. A radar block may override that shape for radar geometry
only. It can select \texttt{wgs-72} or \texttt{wgs-84}, or directly provide
\statevar{reqtr} and one of \statevar{rpolr}, \statevar{flat}, or
\statevar{ecc}. Additional shape values are redundant and may form inconsistent
geometry.

\begin{lstlisting}[style=taosmanual]
*radar 2 station_y alt=112.4 long=65.345 latgd=12.435
                   wgs-72 distn=25.5 diste=15.2 distd=-10.3
\end{lstlisting}

This defines station 2 relative to WGS-72 at 112.4~ft altitude,
$65.345^\circ$~E longitude, and $12.435^\circ$~N latitude, then offsets it
25.5~ft north, 15.2~ft east, and 10.3~ft upward.

\begin{table}[htbp]
  \centering
  \caption{Radar data-block variables.}
  \label{tab:radar-data-block-variables}
  \small
  \begin{tabularx}{0.97\linewidth}{@{}>{\ttfamily}lX>{\ttfamily}lX@{}}
    \toprule
    \normalfont Variable & Description & \normalfont Variable & Description \\
    \midrule
    alt & Radar-station altitude (ft) & reqtr & Equatorial earth radius (ft) \\
    long & Radar-station longitude (deg) & rpolr & Polar earth radius (ft) \\
    latgd & Geodetic radar latitude (deg) & flat & Flattening, $0<\texttt{flat}<1$ \\
    distn & North offset (ft) & ecc & Eccentricity, $0<\texttt{ecc}<1$ \\
    diste & East offset (ft) & wgs-72 & Use WGS-72 earth shape \\
    distd & Down offset (ft) & wgs-84 & Use WGS-84 earth shape \\
    \bottomrule
  \end{tabularx}
\end{table}
